\section{Background and Related Work}
\label{sec:related}

\paragraph{DBpedia.} DBpedia~\cite{auer2007dbpedia, lehmann2015dbpedia} extracts
labels, abstracts, page identifiers, images, categories, redirects, external
links, coordinates and infobox data from every Wikipedia article. Infobox data
is published raw, as properties of the generic \term{dbp:} namespace, and
mapping-based, through community mappings to the hand-built DBpedia Ontology
(\term{dbo:}), which is cleaner but covers less. Resources are named after
their article titles, and language chapters publish their own edition under
their own namespace with links to the English one. Kontokostas et
al.~\cite{kontokostas2012greek} describe what this raises for the Greek
chapter: IRIs in a non-Latin script, language-specific templates and links
across editions. We follow this model for Vietnamese, with a small ontology
where DBpedia lacks the terms we need.

\paragraph{Linked Data.} The Linked Data principles ask for HTTP URIs that return
RDF when looked up and link to other
URIs~\cite{bernerslee2006linked, heath2011linked}. The five-star scheme grades
open data, and Tasks~3 and~4 correspond to its fourth and fifth stars. We serve
the data with \emph{303 See Other} redirects and content
negotiation~\cite{sauermann2008cooluris, hyland2014bestpractices}, describe it
with VoID~\cite{alexander2011void}, and use RDF~1.1~\cite{cyganiak2014rdf} and
SPARQL~1.1~\cite{harris2013sparql}.

\paragraph{Wikidata.} Wikidata~\cite{vrandecic2014wikidata} items carry typed
statements with qualifiers and \emph{sitelinks} to the article about them in
every Wikipedia edition, and the Wikidata Query Service exposes them through
SPARQL~\cite{malyshev2018wdqs}. We use it to select the entities, to obtain
precise dates and quantities, and, through English sitelinks, to link to DBpedia.

\paragraph{Ontology engineering and question answering.} The ontology was
designed from \emph{competency questions}, written before the schema and later
used as tests~\cite{gruninger1995methodology}. DBpedia types and shortcut
relations are derived with the OWL~2~RL profile~\cite{motik2012owl2profiles},
whose rules run by forward chaining in polynomial time, so the derived triples
can be materialised. Retrieval-augmented generation grounds the answer of a
language model in retrieved evidence~\cite{lewis2020rag}; over a knowledge
graph the retrieval can be a SPARQL query written by the model, as on our
question-answering page (Section~\ref{sec:qa}).
